% =============================================================================
% Chapter 1: Introduction (Part I)
% =============================================================================
\mychapter{1. Introduction}

The global adoption of additive manufacturing for the production of structurally critical metal
components has accelerated considerably over the past decade. Industries including aerospace,
automotive, and biomedical engineering have embraced metal AM processes such as laser powder bed
fusion (LPBF), directed energy deposition (DED), and electron beam melting (EBM) for their ability to
produce components with complex internal geometries, reduced material waste, and shortened lead
times compared to traditional machining and casting routes~\citep{Deshpande2024}.

Despite these advantages, metal AM processes are inherently susceptible to the formation of internal
defects during fabrication. These defects arise from the complex physics of rapid solidification,
thermal gradients, powder feedstock variability, and laser or beam parameter instability. The most
commonly observed defect classes include gas porosity arising from entrapped gas in the melt pool,
lack-of-fusion (LoF) voids caused by insufficient energy density and incomplete melting of adjacent
powder particles, micro-cracks resulting from thermal stresses and residual strain accumulation, and
metallic inclusions originating from contaminated powder or redeposited spatter~\citep{Zubayer2023}.
Each of these defect types can initiate fatigue crack growth under cyclic loading, dramatically
reducing component service life and potentially causing catastrophic failure in safety-critical
applications~\citep{Liu2024}.

Effective, reliable, and efficient defect inspection is therefore not merely a quality assurance
concern but a fundamental prerequisite for the safe deployment of AM components. X-ray computed
tomography has established itself as the gold standard for volumetric non-destructive inspection,
capable of imaging internal features at resolutions approaching the micrometre scale without
requiring sectioning or destruction of the component~\citep{Miller2020}. The NIST CoCr AM XCT dataset
used in this work, comprising 900 slices at $1010 \times 980$ pixel resolution with a raw intensity
range of $[0, 59631]$, is representative of the scale and complexity typical of industrial XCT
inspection tasks~\citep{NIST_CoCr}.

However, the volumetric nature of XCT data creates substantial challenges for manual analysis. Expert
radiographers or materials engineers must review vast quantities of image data, a process that is
time-consuming, costly, operator-dependent, and difficult to scale to industrial production volumes.
The primary technical bottleneck in XCT-based defect analysis is not data acquisition but data
interpretation: specifically, the automated detection and localisation of defect regions within
noisy, artefact-affected volumetric images.

Classical approaches based on global thresholding such as Otsu's method~\citep{Otsu1979},
morphological filtering, and edge detection fail under the imaging conditions typical of dense metal
scans, where defects and the surrounding metal matrix occupy similar greyscale intensity ranges and
beam hardening artefacts create systematic intensity gradients across the image
volume~\citep{Bellens2021}. In the NIST AM Bench dataset, Otsu thresholding achieves Dice scores of
only 0.43--0.61 for porosity detection, compared to 0.82--0.91 for CNN-based
approaches~\citep{Wong2022}, confirming the limitations of classical methods for this application.

Deep learning, and in particular fully convolutional architectures such as
U-Net~\citep{Ronneberger2015}, offers a fundamentally different approach. Rather than relying on
hand-crafted intensity thresholds or morphological rules, deep learning models learn rich
hierarchical representations of spatial context, texture, and shape directly from training data. This
enables them to detect and localise defect voxels even under challenging noise and contrast
conditions, and to generalise across different component geometries and scan settings when trained on
sufficiently diverse datasets~\citep{Zubayer2023}.

A critical practical challenge in this thesis is the absence of expert-annotated ground truth masks
for the NIST CoCr AM XCT dataset~\citep{NIST_CoCr}. This is addressed through an automatic label
generation strategy: Otsu thresholding is applied within the metal region only, not the air
background, to generate binary detection labels, which then serve as training supervision for the 2D
U-Net. This approach follows Yosifov~\citep{Yosifov2020}, who used Otsu-derived labels to train a
modified 3D U-Net for XCT analysis of fibre-reinforced polymers, and is justified by the observation
that the trained U-Net already achieves a validation Dice score of 0.68 after 20 epochs, compared to
0.41 for the BCE baseline and 0.38 for the Otsu label generator itself, demonstrating that the model
learns to detect defects beyond the capability of its own training labels.

This thesis addresses the design, implementation, and empirical validation of a deep learning
pipeline for automated XCT defect detection in additively manufactured metal parts. The complete
pipeline, covering preprocessing, automatic label generation, patch extraction, data augmentation,
model training, and evaluation, has been implemented in Python, version-controlled on Git, and
validated on the NIST CoCr \texttt{set1sample2raw} dataset. Part II of this document extends the
validation to real PODFAM process-monitoring data.

\mysection{1.1 Aim}
The aim of this thesis is to develop, implement, and validate an automated deep learning-based
pipeline for the detection of internal defects in XCT scans of additively manufactured metal
components. The work emphasises rigorous empirical methodology: systematic comparison of a deep
learning detector (2D U-Net) against a classical baseline (Otsu thresholding) using standardised
evaluation metrics on the NIST CoCr AM XCT benchmark dataset, followed by assessment of robustness
across preprocessing configurations and loss function formulations.

The ultimate deliverable is a documented, validated, and open-source defect detection pipeline that
can be adopted or extended by manufacturing organisations seeking to automate XCT-based quality
assurance.

\mysection{1.2 Research Questions}
This thesis is structured around the following specific research questions:
\begin{enumerate}
\item \textbf{RQ1}: How effectively can a preprocessing pipeline reduce imaging artefacts and noise
  in XCT scans of AM metal components, as measured by signal-to-noise ratio improvement and visual
  inspection of defect-relevant structures across all four pipeline stages?
\item \textbf{RQ2}: What preprocessing parameter configurations, specifically NLM filter strength,
  BHC polynomial degree, and ring artefact filter radius, produce optimal image quality for
  subsequent automated defect detection?
\item \textbf{RQ3}: How do preprocessing strategies, patch extraction, and data augmentation affect
  the quality of automatically generated detection labels and downstream 2D U-Net detection
  performance under conditions of severe class imbalance (1.2--2.8\% defect voxels)?
\item \textbf{RQ4}: What loss function formulations best address the class imbalance inherent in XCT
  defect data, and how does loss function choice affect detection accuracy as measured by Dice
  coefficient, IoU, precision, and recall?
\item \textbf{RQ5}: Does a 2D U-Net trained on classically-derived pseudo-labels from one industrial
  dataset generalize to a fully held-out sample and to an out-of-distribution sample from the same
  process, as measured by both pixel-level overlap and volume-fraction agreement?
\item \textbf{RQ6}: Are classical thresholding methods (Otsu, Yen, Bernsen) reliable enough to serve
  as an independent, unsupervised comparison baseline across a real industrial sample set, or do they
  carry systematic, sample-dependent failure modes that could be mistaken for genuine measurement?
\end{enumerate}

\mysection{1.3 Contributions}
The main contributions of this thesis are:
\begin{itemize}
\item A complete end-to-end Python pipeline for XCT defect detection covering preprocessing,
  automatic label generation, patch-based data preparation, 2D U-Net training, and evaluation, fully
  version-controlled on Git with MLflow experiment tracking and following FAIR data stewardship
  principles~\citep{Wilkinson2016}.
\item A systematic preprocessing pipeline validated on the NIST CoCr \texttt{set1sample2raw} dataset
  (900 slices, $1010 \times 980$ pixels, raw range $[0, 59631]$), achieving noise standard deviation
  reduction from 20952 to 0.385, a factor of $\times 54{,}500$, and SNR improvement from 1.756 to
  1.773 across the four-stage pipeline.
\item A comprehensive preprocessing parameter sensitivity analysis evaluating NLM filter strength
  ($h \in \{0.5, 1.0, 1.5, 2.0\}$), BHC polynomial degree ($d \in \{1, 2, 3, 5\}$), and ring artefact
  filter radius ($r \in \{5, 10, 15, 25\}$ pixels), with quantitative SNR comparison across all
  configurations.
\item An automatic label generation strategy using two-stage Otsu thresholding with morphological
  cleaning, enabling supervised U-Net training without manual
  annotation~\citep{Yosifov2020,Otsu1979}.
\item Preliminary empirical evidence that the 2D U-Net trained on automatic labels already
  outperforms the Otsu label generator itself (validation Dice 0.68 vs.\ 0.38 after 20 epochs),
  confirming that the model generalises beyond its noisy supervision signal.
\item A documented, reproducible open-source implementation with thesis-ready analysis outputs
  including intensity histograms, SNR plots, parameter tuning comparisons, patch extraction grids,
  and augmentation validation figures.
\item (Part II, answering RQ5) A completed validation of this pipeline on real PODFAM metal AM XCT
  data, including a trained detector evaluated on a fully held-out sample and an additional
  out-of-distribution sample.
\item (Part II, answering RQ6) A documented and independently reproduced classical-thresholding
  calibration failure mode, traced to processed-image mean intensity and confirmed on two separate
  PODFAM samples rather than dismissed as a one-off anomaly.
\end{itemize}
